% Transcription — fonds Alexandre Grothendieck, shelfmark 62, batch 2 (pages 21-40).
% Established from the facsimile archives/batches/62/batch-02.pdf (a copy of
% 62.pdf fetched through the project relay, no cover sheet: PDF page k =
% archivists' page 20+k), per the skill transcribe-grothendieck. The folder is
% ninety-six pages in five batches; this is the second.
%
% Pass: Opus 5.5 (claude-opus-5-5), 2026-09-26 — first pass, unchecked against the pages by a human.
%
% Pages skipped: 25, 28, 30, 32, 34, 36, 38 and 40, typed leaves (the other
% sides of the sheets he wrote on) with nothing of his mathematics: pages of
% two Bourbaki-style typescripts, « n° 276 » (filtered modules, m-adic
% topologies; typed pp. 11 and 18) and « n° 272 » (valuations; typed pp. 39,
% 43-47), some upside down; page 25 carries only an ink doodle. Page 26 is a
% blank sheet. Pages 21 and 23 are typed pages 9 and 10 of the same « n° 276 »
% that carry pencil marginalia; by the settled rule (a typed verso is kept
% when it carries something in his hand) they are kept, with one \note each
% on the typed content and the marginalia given as far as they read. The
% attribution of that pencil to him is not certain, and the \note says so.
%
% Pages summarised: none. Pages 22 and 24 are scratch sheets of computation
% (no prose), given formula by formula in page order.
%
% His pagination: none on these leaves. Numbered runs of his own: the
% numbered formulas 1)-8) (pp. 27-29), (9) (p. 35), (10), (11), (12) (p. 39),
% one continuous run across the odd pages 27-39; a bold « 5 » at the end of a
% line on p. 27 is recorded in a \note. No dates.
%
% What the batch is about. Pages 27-39 are a connected draft on the Legendre
% cross-ratio parameter and j: P^1_S with the three sections 0, 1, infinity;
% its automorphism group S_3 (phi, psi of order 3, the transpositions lambda,
% mu, nu, their multiplication table and conjugations), an extra involution
% sigma(t) = (t-2)/(2t-1) commuting with S_3 (the normaliser S_3 x Z/2 in the
% margin of p. 29); fixed points (1,-1), (infinity,1/2), (0,2) and the roots
% of Phi_6; the quotients P^1_S -> P^1_S/(Z/3) -> P^1_S/S_3 away from
% characteristics 2 and 3 (exact sequences with a twisted G_m, the
% ramification picture on p. 35, S' = S(sqrt(-3))), the explicit maps
% g(s) = s^2 and f(t) = (1/3 sqrt(-3)) (t+1)(2t-1)(t-2)/(t(t-1)) (p. 37, via
% conjugation by u(t) = (t-x')/(t-x'') and h(t) = -t^3), whence
% j = -(1/27) ((t+1)(2t-1)(t-2)/(t(t-1)))^2 as a Galois S_3-cover (p. 39).
% Pages 22 and 24 carry the matching computations in the roots x_1, x_2, x_3
% of a cubic: j = (27c + 2a^3 - 9ab)^2 / (-27 Delta), the cross-ratio
% t = (x_3-x_1)/(x_3-x_2), and y^2 = 4x^3 - h(x+1), h = 27j/(j - 2^6 3^3).
%
% Notation fixed once for the batch, following batch 1: projective line as
% \mathbf{P}^1_S, the symmetric group as \mathfrak{S}_3, the multiplicative
% group as \mathbf{G}_m, integers as \mathbf{Z}; the third transposition,
% written like a \vartheta, is read \nu; left exponents (conjugation) as
% {}^{\varphi}\lambda. Inside formulas an unreadable index or exponent is
% written « ? », as in batch 1.
%
% Hardest pages: 29, 31, 33 (fast connective prose between the formulas),
% 37 (the opening paragraph, mostly \ill), 39. Readings to check first: the
% index of \mathfrak{S}_6 and the word before it at the foot of p. 27; the
% word « normalisateur » in the margin of p. 29; « (\mathbf{G}_a)_S » on
% p. 39; the attribution of the pencil on pp. 21 and 23.

\documentclass[11pt,a4paper]{article}
% Le préambule commun à toutes les transcriptions du fonds Grothendieck.
%
% Il tient en une idée : **l'incertitude se compose, elle ne se résout pas.**
% Une transcription qui lisse un mot douteux a détruit la seule chose pour
% laquelle on la faisait — un lecteur doit pouvoir savoir, à chaque endroit,
% ce qui était sur le feuillet et ce qui a été ajouté. Les six macros qui
% suivent sont donc l'appareil critique, et non de la décoration.
%
% Les mêmes commandes sont reconnues par `scripts/render.mjs`, qui produit la
% vue de lecture du site. Une macro ajoutée ici doit l'être aussi là-bas,
% faute de quoi le rendu échouera bruyamment — ce qui est voulu : un rendu
% silencieusement faux est pire qu'un rendu absent.

\ProvidesPackage{grothendieck}[2026/08/08 Transcriptions du fonds Grothendieck]

\usepackage{amsmath,amssymb,amsthm}
\usepackage{mathrsfs}
\usepackage{stmaryrd}
% Les diagrammes commutatifs sont partout dans ce fonds : c'est la langue
% même de la géométrie algébrique de Grothendieck, et une transcription qui
% les rendrait en prose ne transcrirait rien.
\usepackage{tikz-cd}
% Français par défaut — c'est la langue des feuillets — mais l'anglais est
% chargé aussi : la traduction et le résumé sont des documents anglais, et la
% typographie française (espaces avant les deux-points) y serait une faute.
% Ces documents basculent par \selectlanguage{english} en tête de corps.
\usepackage[english,french]{babel}
\usepackage{fontspec}
\usepackage{geometry}
\geometry{a4paper,margin=25mm,marginparwidth=28mm}
\usepackage{marginnote}
\usepackage{xcolor}
% ulem, not soul. `soul`'s \st and \ul are built on a character-by-character
% scan that XeLaTeX's font machinery breaks: the document fails with an
% undefined control sequence rather than degrading. ulem does the same two jobs
% and is Unicode-engine safe. `normalem` stops it hijacking \emph, which the
% transcriptions use for Grothendieck's own emphasis.
\usepackage[normalem]{ulem}
\usepackage{hyperref}
\hypersetup{colorlinks=true,linkcolor=black,urlcolor=black,pdfborder={0 0 0}}

% --- Métadonnées du lot -----------------------------------------------------
% Reprises telles quelles par le script de rendu, qui les affiche en tête de
% la vue de lecture. Elles fixent ce que le fichier prétend couvrir : sans
% elles, une transcription flotte sans point d'attache dans un fonds de seize
% mille feuillets.

\newcommand{\folder}[1]{\gdef\grothFolder{#1}}
\newcommand{\batch}[1]{\gdef\grothBatch{#1}}
\newcommand{\leaves}[2]{\gdef\grothFirst{#1}\gdef\grothLast{#2}}
\newcommand{\foldertitle}[1]{\gdef\grothFoldertitle{#1}}
\gdef\grothFolder{?}\gdef\grothBatch{?}\gdef\grothFirst{?}\gdef\grothLast{?}\gdef\grothFoldertitle{}

% --- Le résumé --------------------------------------------------------------
%
% Il ouvre la lecture modernisée, sous le titre « Résumé », et il est composé
% en petit. Ce n'est pas un ornement : c'est le seul passage du document qui
% ne soit pas des mathématiques, et un lecteur doit pouvoir le reconnaître
% avant de l'avoir lu — puis le sauter s'il connaît déjà le sujet, ou s'y
% arrêter s'il ne le connaît pas. Le filet qui le suit marque la reprise du
% corps.

\newenvironment{resume}
  {\small\setlength{\parskip}{0.45em}}
  {\par\medskip\hrule\medskip}

% --- Mots-clés ---------------------------------------------------------------
%
% En anglais, en clôture du résumé de la lecture modernisée : le vocabulaire
% moderne sous lequel chercher ce que les feuillets construisent. Le manifeste
% du site (`npm run manifest`) les extrait pour étiqueter la cote entière —
% cette ligne est la source unique des tags, il n'y a pas de fichier à côté.
\newcommand{\keywords}[1]{\par\smallskip\noindent
  {\footnotesize\textsc{Keywords} --- \textit{#1}}\par}

% --- L'appareil critique ----------------------------------------------------

% Le numéro de feuillet, en marge. C'est l'ancre qui permet de régler un
% désaccord sur une formule en regardant *un* feuillet plutôt qu'une chemise
% de six cents. Le site s'en sert aussi pour tourner les pages du fac-similé
% quand on fait défiler la transcription.
\newcommand{\leaf}[1]{%
  \marginnote{\footnotesize\color{gray!70}\textsf{#1}}%
  \label{leaf:#1}\ignorespaces}

% Illisible : on ne devine pas. Un mot inventé qui se lit comme les autres est
% le pire résultat possible.
\newcommand{\ill}{\textcolor{red!60!black}{[\,\ldots\,]}}

% Lecture douteuse : le mot est donné, et signalé comme incertain.
\newcommand{\uncertain}[1]{\uline{#1}}

% Ajout de l'éditeur — une lettre manquante, un mot sous-entendu.
\newcommand{\add}[1]{\textcolor{blue!50!black}{[#1]}}

% Biffé par Grothendieck lui-même. Ce qu'il a rayé fait partie du document :
% une rature dit souvent où la pensée a changé de direction.
\newcommand{\struck}[1]{\sout{#1}}

% Note du transcripteur, sur l'état matériel du feuillet.
\newcommand{\note}[1]{\par\smallskip{\small\color{gray!60!black}\textit{[#1]}}\par\smallskip}

% Note marginale de Grothendieck, distincte du corps du texte.
\newcommand{\marginal}[1]{\par\smallskip\noindent\hspace{1em}%
  {\small\color{gray!60!black}#1}\par\smallskip}

% --- Titre ------------------------------------------------------------------

\AtBeginDocument{%
  \begin{center}
    {\large\bfseries Fonds Alexandre Grothendieck --- cote n\textsuperscript{o}~\grothFolder}\\[2pt]
    {\itshape\grothFoldertitle}\\[4pt]
    {\small feuillets \grothFirst--\grothLast\ (lot \grothBatch)}\\[2pt]
    {\footnotesize\color{gray!60!black}Transcription établie d'après le fac-similé numérisé
     par l'Université de Montpellier.\\
     Les lectures incertaines sont soulignées, les illisibles notées [\,\ldots\,],
     les ajouts entre crochets.}
  \end{center}
  \vspace{1em}}

\endinput


\folder{62}
\batch{2}
\pages{21}{40}
\dating{[à partir de 1963]}
\watermark{Édition de démonstration}
\foldertitle{Courbes elliptiques (généralités - vieille rédaction) : notes manuscrites (s.d.)}

\begin{document}

\page{21}
\note{page 9 d'un tapuscrit dactylographié « n° 276 » (modules filtrés
complets : Proposition 4, un module filtré est complet ssi toute série dont le
terme général tend vers 0 converge ; Corollaires 1 et 2), sans rapport direct
avec les courbes elliptiques. Il porte au crayon, dans la marge gauche, en
face du Corollaire 1, une remarque d'une main qui paraît être la sienne, sans
certitude ; un trait ondulé souligne « résulte »}
\marginal{fait \uncertain{général} \ill{} \ill{} \ill{} \uncertain{métrisables} !}

\page{22}
\note{feuille de calcul, à l'encre, sans rédaction : le $j$ d'une cubique en
fonction de ses racines $x_1,x_2,x_3$. Les calculs sont donnés dans l'ordre de
la page}
\[
\frac{(t+1)(2t-1)(t-2)}{t(t-1)} =
\frac{(2x_1-x_2-x_3)(2x_2-x_3-x_1)(2x_3-x_1-x_2)}{(x_1-x_2)(x_2-x_3)(x_3-x_1)} .
\]
\note{entre $2x_1$ et $x_2$, au numérateur, un signe surchargé}
\struck{$x_1+x_2+x_3$}
\[
\boxed{\ j = -\frac{1}{27\Delta}\bigl[(2x_1-x_2-x_3)(2x_2-x_3-x_1)(2x_3-x_1-x_2)\bigr]^2\ }
\]
\note{la formule biffée, à gauche du cadre, est lue sous réserve}
\[
\sum x_i = -a, \qquad \sum x_ix_j = b, \qquad x_1x_2x_3 = -c
\qquad\qquad
\Delta = \bigl((x_1-x_2)(x_2-x_3)(x_3-x_1)\bigr)^2
\]
\note{le signe devant $c$ est écrit sur une tache}
\[
j = \frac{\bigl[(3x_1+a)(3x_2+a)(3x_3+a)\bigr]^2}{-27\,\Delta}
\]
\struck{$\dfrac{27}{-\Delta}\Bigl[\bigl(\tfrac{a}{3}+x_1\bigr)\bigl(\tfrac{a}{3}+x_2\bigr)\bigl(\tfrac{a}{3}+x_3\bigr)\Bigr]^2$}
\note{sous cette ligne biffée, un $\tfrac{a}{3}$ entouré et surchargé}
\[
\begin{aligned}
\sqrt{-27\Delta j} &= a^3 + 3a^2(x_1+x_2+x_3) + 9a(x_1x_2+x_2x_3+x_3x_1) + 27\,x_1x_2x_3\\
&= a^3 - 3a^3 + 9ab - 27c\\
&= -2a^3 + 9ab - 27c
\end{aligned}
\]
\note{dans la première ligne, le signe après $a^3$, l'exposant de $3a^2$ et le
coefficient $9$ sont écrits sur d'autres signes ; une tache avant
$x_1x_2x_3$}
\[
\boxed{\ j = \frac{(27c + 2a^3 - 9ab)^2}{-27\,\Delta}\ }
\]
\struck{$\prod(x_1^2+x_2^2-2x_1x_2) =$}

\page{23}
\note{page 10 du même tapuscrit « n° 276 » (fin de la preuve du Corollaire 2 ;
Proposition 5 : si $A$ est séparé, resp.\ complet, pour la topologie
$\mathfrak{m}$-adique, l'anneau $A[[X]]$ filtré par les puissances de l'idéal
engendré par $\mathfrak{m}$ et $X$ est séparé, resp.\ complet). Au crayon, de
la même main qu'à la page 21 : le mot « \uncertain{engendré} » est écrit
au-dessus de « par $\mathfrak{m}$ et $X$ » dans l'énoncé, et une note court en
biais dans la marge gauche}
\marginal{Dire \ill{} \uncertain{topologie} \ill{} \ill{} \uncertain{celle des}
\ill{} \ill{} \ill{} \uncertain{coefficients}}

\page{24}
\note{feuille de calcul, à l'encre : forme normale, birapport et racines de la
cubique}
\[
y^2 = 4x^3 - h(x+1) \qquad\qquad h = 27j\,(j - 2^6 3^3)^{-1}
\]
\struck{$y^2 = ax^3 + bx^2 + cx + d$}
\[
y^2 = x^3 + ax^2 + bx + c
\]
\note{sous $a$, $b$, $c$ sont écrits les chiffres $1$, $2$, $3$ ; le $a$ est
écrit sur une autre lettre}
\[
\begin{array}{cccc}
x_1 & x_2 & x_3 & x_4\\
\downarrow & \downarrow & \downarrow & \downarrow\\
\infty & 0 & 1 & t
\end{array}
\]
\note{le $x_4$ est surchargé ; à droite, un petit croquis des quatre points
$\infty, 0, 1, t$ sur une droite, reliés deux à deux par des crochets}
\[
((0,1,\infty,t)) = ((x_1,x_2,x_3,x_4))
\]
\note{une accolade sous « $0,1$ »}
\[
\frac{x_1-x_3}{x_1-x_4} : \frac{x_2-x_3}{x_2-x_4}
= \frac{0-\infty}{0-t} : \frac{1-\infty}{1-t}
= \frac{1-t}{-t} = -\frac{1}{t} + 1
\]
\[
\frac{t-1}{t} \qquad\qquad \frac{t-\infty}{t-0} : \frac{1-\infty}{1-0}
\]
\[
t = \frac{1-\infty}{1-0} : \frac{t-\infty}{t-0}
= \frac{x_3-x_1}{x_3-x_2} : \frac{x_4-x_1}{x_4-x_2}
\]
\uncertain{Si} $x_4=\infty$, \ill{} \ill{} \ldots
\[
\boxed{\ t = \frac{x_3-x_1}{x_3-x_2}\ }
\qquad
t-1 = \frac{x_2-x_1}{x_3-x_2}
\qquad
\frac{c-a}{c-b} \;\Big|\; \frac{c-b}{c-a}
\]
\[
\frac{a-b}{a-c} \qquad \frac{a-b}{c-b} \cdot \frac{b-a}{b-c}
\]
\note{ces fractions en $a,b,c$ sont griffonnées au bord droit, la dernière
coupée par le bord ; le signe entre les deux dernières est une tache}
\[
t+1 = \frac{2x_3-x_1-x_2}{x_3-x_2}
\qquad
2t-1 = \frac{x_3-2x_1+x_2}{x_3-x_2}
\qquad
t-2 = \frac{-x_3-x_1+2x_2}{x_3-x_2}
\]
\note{dans la dernière fraction, le signe devant $x_1$ est surchargé}

\page{27}
$\mathbf{P}^1_S$ \uncertain{sections} $0,1,\infty$. Le groupe des automorphismes
s'identifie au groupe $\mathfrak{S}_3$ des permutations des trois points
$0,1,\infty$, groupe \uncertain{qui a} $6$ éléments
$1,\varphi,\psi,\lambda,\mu,\nu$
\note{la dernière lettre, écrite comme un $\vartheta$, est lue $\nu$ ici et
dans toute la suite ; un long trait vertical double et une boucle, à
l'encre brune, dans la marge gauche en haut de la page}
\[
(1)\quad
\begin{cases}
\varphi(t) = 1 - \dfrac{1}{t}\\[2ex]
\psi(t) = \dfrac{1}{1-t}
\end{cases}
\]
$2$ permutations circulaires non triviales sur $0,1,\infty$
\[
(2)\quad
\begin{cases}
\lambda t = \dfrac{1}{t}\\[1ex]
\mu t = 1-t\\[1ex]
\nu t = 1 + \dfrac{1}{t-1} = \dfrac{t}{t-1}
\end{cases}
\qquad\qquad
\begin{array}{c} 1\\ \infty\\ 0 \end{array}
\]
transposition de $0,\infty$ ; transposition de $0,1$ ; transposition de
$1,\infty$ (dans l'ordre des trois lignes).
\note{à droite des trois formules, séparés par un double trait vertical, les
points $1$, $\infty$, $0$, chacun en face de sa ligne : le point fixé par la
transposition}
\marginal{\uncertain{renommer} : \uncertain{$\lambda t = 1-t$}, \uncertain{$\mu t
= t/(t-1)$}, \uncertain{$\nu t = 1/t$} ??}
\note{note en biais au bord droit, délimitée par une accolade, en partie
coupée par le bord}
\[
3)\ \text{a)}\ \boxed{\begin{array}{ll} \varphi^2 = \psi & \psi^2 = \varphi\\
& \varphi\psi = \psi\varphi = 1 \end{array}}
\quad
\text{b)}\ \boxed{\lambda^2 = \mu^2 = \nu^2 = 1}
\quad
\text{c)}\ \boxed{\begin{array}{l} \lambda\mu = \mu\nu = \nu\lambda = \varphi\\
\mu\lambda = \nu\mu = \lambda\nu = \psi \end{array}}
\]
\[
\text{d)}\ \boxed{\begin{array}{l|l}
\varphi\lambda = \mu\varphi = \nu & \psi\nu = \mu\psi = \lambda\\
\varphi\mu = \nu\varphi = \lambda & \psi\mu = \lambda\psi = \nu\\
\varphi\nu = \lambda\varphi = \mu & \psi\lambda = \nu\psi = \mu
\end{array}}
\]
\note{dans le cadre a), sous $\varphi^2=\psi$, un mot biffé et illisible ; le
cadre d) est au bord droit de la page}

$1,\varphi,\psi$ \ill{} \ill{} \uncertain{cyclique} d'ordre $3$,
\uncertain{engendré par} $\varphi$ et $\psi$, invariant dans $\mathfrak{S}_3$.
Le quotient \uncertain{est} $\mathbf{Z}/2$ et $\mathfrak{S}_3$ est
\uncertain{produit semi-direct} ; il y a trois sous-groupes \ill{},
\uncertain{engendrés resp.\ par} $\lambda,\mu,\nu$ (les \uncertain{seuls
éléments} d'ordre $2$)
\[
4)\quad
\begin{cases}
{}^{\varphi}\lambda = \mu \quad {}^{\varphi}\mu = \nu \quad {}^{\varphi}\nu = \lambda\\
{}^{\psi}\nu = \mu \quad {}^{\psi}\mu = \lambda \quad {}^{\psi}\lambda = \nu
\end{cases}
\qquad
5)\quad
\begin{cases}
{}^{\lambda}\varphi = {}^{\mu}\varphi = {}^{\nu}\varphi = \psi\\
{}^{\lambda}\psi = {}^{\mu}\psi = {}^{\nu}\psi = \varphi
\end{cases}
\]
\note{4) et 5) sont écrits à droite du paragraphe qui précède ; l'exposant à
gauche note la conjugaison}

De plus, on a \add{\uncertain{de} $\mathbf{P}^1$} un automorphisme
\struck{\ill{}} de carré $1$, donné par
\[
6)\quad \sigma t = \frac{t-2}{2t-1} \qquad [\sigma^2 = \mathrm{id}]
\]
\note{un « $5$ » appuyé termine la ligne « De plus, on a \ldots\ automorphisme » ;
son rôle n'est pas clair. La lettre $\sigma$ est écrite sur une autre, les
deux fois}
\uncertain{Normalisateur} : $\mathfrak{S}_6$ (c'est d'ailleurs le
\uncertain{seul}, \ill{}
\note{l'indice de $\mathfrak{S}$, lu $6$, est peu net ; la phrase se poursuit à
la page 29}

\page{29}
l'identité). Il \ill{} $0,1,\infty$ en
\[
7)\quad
\begin{cases}
\sigma 0 = 2\\
\sigma 1 = -1\\
\sigma\infty = \tfrac12
\end{cases}
\]
\marginal{\uncertain{Prouver} : le \uncertain{normalisateur} de $\mathfrak{S}_3$
dans $\operatorname{Aut}\mathbf{P}^1_S$ est
$\mathfrak{S}_3\times(\mathbf{Z}/2)$}
\note{note en biais dans la marge gauche ; le $\mathbf{Z}$ de
$\mathbf{Z}/2$ est surchargé et entouré}
La différente de $\mathfrak{S}_3$ \ill{} \ill{} \uncertain{pts} de
ramification --- \ill{} formé\supplied{e} de l'ens.\ des pts \uncertain{fixes}
de \uncertain{l'un} des cinq \ill{} $\lambda\mu\nu\varphi\psi$ \uncertain{distincts}
de l'identité. Pour $\lambda,\mu,\nu$, on trouve les sections $(1,-1)$
$(\infty,\tfrac12)$ $(0,2)$, [\uncertain{stables} par $\sigma$]. Pour
$\varphi$ (et $\psi$, c'est pareil), on trouve
\note{le couple $(1,-1)$ est écrit sur un autre, en surcharge appuyée}
\[
1 - \frac1t = t \quad \ldots
\]
\[
8)\qquad t^2 - t + 1 = 0
\]
i.e.\ l'équation cyclotomique $\boxed{\Phi_6(t) = 0}$. \struck{$6$}

\struck{\ill{}} \uncertain{C'est} aussi \add{$\overline{S}$} l'un des pts
fixes sous $\sigma$.
\note{le $\overline{S}$, surligné, est écrit au-dessus de « l'un » ; sa place
dans la phrase n'est pas sûre}
[Il est non ramifié \add{sur $S$} aux pts $s\in S$ \uncertain{tels que} $k(s)$
\uncertain{soit} de car.\ $\neq 2,3$]. $\mathbf{Z}_2 =
\mathfrak{S}_3/(\mathbf{Z}/3)$ y \uncertain{opère}, de façon non triviale aux
pts \ill{} \ill{} \ill{} non ramifié sur $S$.

D'autre part, le \uncertain{schéma} \add{$G$} des automorphismes de
$\mathbf{P}^1$ qui laissent $\overline{S}$ \struck{\ill{}} invariant est une
extension de $\mathbf{Z}_2$ [automorphismes de $\overline{S}$] par
$(\mathbf{G}_m)_S$ \emph{tordu} ---- [par le rev.\ cyclotomique]
\marginal{les \uncertain{pts} \ill{} \ill{} \ill{} \ill{} (car.\ $2$ ou $3$) de $S$}
\note{note dans la marge gauche, séparée du texte par un trait vertical, en
face de ce paragraphe}

\page{31}
\begin{tikzcd}
e \arrow[r] & \mathbf{G}_m^{\mathrm{tordu}} \arrow[r] & G \arrow[r] & \mathbf{Z}/2 \arrow[r] & e \\
e \arrow[r] & \mathbf{Z}/3 \arrow[r] \arrow[u] & \mathfrak{S}_3 \arrow[r] & \mathbf{Z}/2 \arrow[r] \arrow[u, no head, "\Vert" description] & e
\end{tikzcd}

\note{sur la page, une seule flèche verticale, $\mathbf{Z}/3\to\mathbf{G}_m$,
et un signe $\|$ entre les deux $\mathbf{Z}/2$ ; aucune flèche n'est tracée de
$\mathfrak{S}_3$ vers $G$. « tordu » est écrit en exposant de $\mathbf{G}_m$}

On \uncertain{supposera} \uncertain{désormais} qu'il n'y a pas de pts $s$ de
car.\ $2$ ou $3$ dans $S$. Alors \struck{\ill{} opère} $\mathfrak{S}_3\times
\mathbf{Z}_2$ opère « tamely » sur $\mathbf{P}^1_S$ [en effet, il est compatible
avec la fibration $\mathbf{P}^1_S\to S$, et opère « tamely » sur les fibres
\ldots] le quotient \uncertain{s'identifie} \add{\ill{} par le groupe ou un
ss-groupe \ill{} pour \ill{}} \ill{} à $\mathbf{P}^1_S$ [car il est \emph{lisse}
sur $S$, et rationnel sur $S$ d'après Lüroth]. Les sections de pts fixes
$(0,1,\infty)$ $(2,-1,\tfrac12)$ $(\overline{S}=x',x'')$ --- qui sont
disjointes, et chacune stable sous $\mathfrak{S}_3$, \add{\uncertain{chacun de
ces ensembles}} sur lequel \struck{il} opère de façon \struck{\ill{}}
transitive sans \ill{} \ill{} chaque fibre ; donnent \ill{} \ill{} \ill{} de
$\mathbf{P}^1_S/\mathfrak{S}_3$ trois \emph{sections} distinctes [à
expliciter\ldots] qu'on peut utiliser pour identifier
$\mathbf{P}^1_S/\mathfrak{S}_3$ à $\mathbf{P}^1_S$ \uncertain{muni} des
sections $\infty,0,1$. De \uncertain{même}, \struck{de se donner des
\ill{}} $\mathbf{P}^1_S/(\mathbf{Z}/3)$ est muni d'abord de deux sections,
qui correspondent à $(0,1,\infty)$ et $(2,-1,\tfrac12)$, et qui
\note{« s'identifie » est traversé par le trait qui appelle l'ajout en
interligne ; le premier des deux $2$ de $(2,-1,\tfrac12)$ est surchargé}

\page{33}
sont les sections de ramification de $\mathbf{P}^1_S$ sur
$\mathbf{P}^1_S/(\mathbf{Z}/3)$. Prenons \ill{} \ill{} \ill{}
$S' = S(\sqrt{-3})$ de $S$, d'où dans $\mathbf{P}^1_{S'}$ \add{deux sections}
\struck{un diviseur \ill{}} \struck{\ill{}} par $\overline{S}'$, soient $x'$ et
$x''$, qui donnent sur $\mathbf{P}^1_{S'}/(\mathbf{Z}/3)$ deux sections ;
\uncertain{prenant} pour l'une comme \struck{la} section \uncertain{unité}
($\infty$ et $0$ étant déjà \uncertain{choisis} ci-dessus) \ill{} \ill{} que
l'autre est $-1$. Cela définit un isomorphisme
\[
(\mathbf{P}^1_S/(\mathbf{Z}/3))\times_S S' \simeq \mathbf{P}^1_{S'} .
\]
Ainsi, $\mathbf{P}^1_S/(\mathbf{Z}/3)$ est can.\ isomorphe à un fibré principal
de groupe $\mathbf{G}_{m,S}$, [associé à un faisceau inversible
$\Lambda_S$], $\Lambda_S$ étant muni d'une bi-section \add{non ramifiée}
\struck{\uncertain{symétrique}} non ramifiée, ou (ce qui revient au même)
d'une forme \uncertain{quadratique} non \uncertain{dégénérée} ; le fait de
munir de cette bisection \ill{} \ill{} autre que celle du revêtement
$S' = S(\sqrt{-3})$.

\uncertain{Posant} $S' = S(\sqrt{-3})$, on obtient la \uncertain{forme}
suivante
\note{le reste de la page est blanc, sauf un mot isolé, « \uncertain{suivante} »
ou « \uncertain{suivants} », en tête de la dernière ligne}

\page{35}
\[
\begin{array}{lllcccc}
t & \mathbf{P}^1 & & \{0,1,\infty\} & \{2,-1,\tfrac12\} & x' = \tfrac12(1+\sqrt{-3}) & x'' = \tfrac12(1-\sqrt{-3})\\
\mathbf{Z}/3 & \big\downarrow & f & \big| & \big| & \big\| & \big\|\\
s & \mathbf{P}^1 & & \infty & 0 & 1 & -1\\
\mathbf{Z}/2 & \big\downarrow & g & \big\| & \big\| & \big\downarrow & \swarrow\\
r & \mathbf{P}^1 & & \infty & 0 & 1 &
\end{array}
\]
\note{schéma de la page, redessiné en tableau. Trois droites
$\mathbf{P}^1$ de coordonnées $t$, $s$, $r$, reliées par $f$ (groupe
$\mathbf{Z}/3$) et $g$ (groupe $\mathbf{Z}/2$). Sur la page : au-dessus de
$s=\infty$ et $s=0$, les fibres $\{0,1,\infty\}$ et $\{2,-1,\tfrac12\}$,
marquées « non ramifié » ; au-dessus de $s=1$ et $s=-1$, les points $x'$ et
$x''$, marqués « ramifié » (traits doubles) ; au-dessus de $r=\infty$ et
$r=0$, les points $s=\infty$ et $s=0$, marqués « ramifié » ; $s=1$ et $s=-1$,
réunis par une accolade, vont tous deux sur $r=1$, marqué « non ramifié ».
Les deux fibres du haut sont écrites au-dessus de chiffres fortement
surchargés ; l'étiquette de $f$ est biffée ; un signe biffé après le $1$ de
la ligne $s$}

Pour déterminer explicitement $f$ et $g$, on procède ainsi. [On
\struck{\ill{}} peut supposer que \add{$S = \operatorname{Spec}(A)$,
$A = \mathbf{Z}[2^{-1},3^{-1}]$} \struck{$S = \operatorname{Spec}(\mathbf{Z}[2^{-1},3^{-1}])$}]
[\ill{} \ill{} \ill{} \ill{} \ill{} \ill{} \ill{} \ill{} \ill{} \ill{}
$\operatorname{Spec}(A)$]. \ill{} \ill{} \ill{} \ill{} \ill{} \ill{}
restriction de $g$ : $\mathbf{P}^1_S-(\infty)-(0)$ est \ill{}
\struck{\ill{}} \ill{} endomorphisme de $\mathbf{G}_m$ dans lui-même, donc est
de la forme $t\mapsto\lambda t^n$, où $n\in\mathbf{Z}$, et
$\lambda\in H^0(S,\underline{O}_S^*)$. Comme $g(\infty)=\infty$, on a $n>0$,
et comme \ill{} $S$ \ill{} \add{non ram.} \add{(d'ordre $2$,} on a $n=2$.
D'autre part \add{$g(1)=1$, d'où $\lambda=1$}
\struck{$\lambda$ \ill{} \ill{} \ill{} $2^\alpha 3^\beta$. \ill{} \ill{} $A$,
donc \ill{} de la forme}
\note{l'ajout « $g(1)=1$, d'où $\lambda=1$ » est écrit sur la fin de la ligne,
qui est barrée d'un double trait, et coupé par le bord droit ; la ligne
suivante est biffée en entier}
\[
(9)\qquad r = g(s) = s^2
\]
\emph{Cette formule est valable \uncertain{sur tout} le $\mathbf{P}^1$} [\ill{}
\ill{} \ill{} \ill{}, et justifier cette notation ---]

\page{37}
\ill{} \ill{} \ill{} \ill{} \ill{} $f$, \ill{} \ill{} \add{\uncertain{comme
dessus}} \add{\ill{}} \ill{} \ill{} \ill{} \ill{} \ill{} \ill{} \ill{} $f\colon
\mathbf{P}^1_{S'}\to\mathbf{P}^1_{S'}$ \ill{} \ill{} \ill{} \ill{} \ill{}
\ill{} \ill{} ; et \ill{} \ill{} \ill{} \ill{} \ill{} section (p.\ ex.\ $0$)
\uncertain{disjointe} de la ramification --- \ill{} \ill{} \ill{} sections
\uncertain{disjointes} --- \ill{} \ill{}. \ill{} \ill{}, \ill{} \ill{}
\ill{} ramification : \ill{} en $\infty$ et \ill{} --- \ill{}
\note{paragraphe écrit vite, lu seulement par fragments ; la formule
$f\colon\mathbf{P}^1_{S'}\to\mathbf{P}^1_{S'}$ est soulignée et en partie
barrée}
\[
s = f(t) = \frac{1 + \Bigl(\dfrac{t-x'}{t-x''}\Bigr)^3}{1 - \Bigl(\dfrac{t-x'}{t-x''}\Bigr)^3}
= \frac{(t-x'')^3 + (t-x')^3}{(t-x'')^3 - (t-x')^3}
\]
\[
= \frac{2t^3 - 3t^2(x''+x') + 3t(x'^2+x''^2) - (x'^3+x''^3)}
{3t^2(x'-x'') + 3t(x''^2-x'^2) - (x''^3-x'^3)}
\]
\note{dans les deux fractions intérieures, un signe biffé devant chaque $t$ ;
les signes « $-$ » des deux dernières lignes sont écrits sur des taches. Un
long trait sinueux sépare la page en deux colonnes : à gauche le diagramme et
le calcul de $f$ comme composé, à droite le calcul de $x'$, $x''$}

\begin{tikzcd}
\mathbf{P}^1 \arrow[r, "u"] \arrow[d, "f"'] & \mathbf{P}^1 \arrow[d, "h"] \\
\mathbf{P}^1 \arrow[r, "v"'] & \mathbf{P}^1
\end{tikzcd}

\note{autour du diagramme : en haut à gauche, « $\infty$ » au-dessus de
« $(x',x'')$ », et $u(t) = \dfrac{t-x'}{t-x''}$ ; en haut à droite,
« $1$ » au-dessus de « $(0,\infty)$ » ; le long de la flèche de droite,
$h^*(t) = -t^3$ ; en bas, $(1,-1)$ et $\infty$ sous le $\mathbf{P}^1$ de
gauche, $v(t) = \dfrac{1-t}{1+t}$, et $(0,\infty)$ et $-1$ sous celui de
droite ; plus bas, une flèche vers la gauche marquée
$v^{-1}(t) = \dfrac{1-t}{1+t}$. La flèche $u$ a une pointe appuyée à son
extrémité gauche aussi}
\[
f = v^{-1}hu
\]
\[
\begin{aligned}
f^*(t) &= u^*h^*v^{-1*}(t) = u^*h^*\Bigl(\frac{1-t}{1+t}\Bigr)\\
&= u^*\,\frac{1-h^*t}{1+h^*t} = u^*\,\frac{1+t^3}{1-t^3}\\
&= \frac{1+u^*(t)^3}{1-u^*(t)^3} =
\end{aligned}
\]

Or ($x'$ et $x''$ \uncertain{étant racines} \emph{primitives}
\uncertain{6ièmes} de $1$)
\[
x'^3 = x''^3 = -1,
\]
d'où
\[
x''^3 - x'^3 = 0
\]
\[
\begin{cases} x'+x'' = 1\\ x'x'' = 1 \end{cases}
\]
\[
\begin{aligned}
x'^2 + x''^2 &= (x'+x'')^2 - 2x'x'' = -1\\
(x'-x'')^2 &= (x'+x'')^2 - 4x'x'' = -3\\
x'-x'' &= \sqrt{-3}\\
x''^2 - x'^2 &= (x''-x')(x''+x') = -\sqrt{-3}
\end{aligned}
\]
\struck{$x''^3 - x'^3 = (x''-x')(x''^2+x'x''+x'^2) = -\sqrt{-3}($}
\[
f(t) = \frac{1}{3\sqrt{-3}}\,\frac{2t^3 - 3t^2 - 3t + 2}{t(t-1)}
= \frac{1}{3\sqrt{-3}}\,\frac{(t+1)(2t-1)(t-2)}{t(t-1)}
\]
\note{le signe devant $3t$ est écrit sur une tache ; le quotient
$(t+1)(2t-1)(t-2)/t(t-1)$ est celui de la première ligne de la page 22}

\page{39}
On obtient ainsi
\[
(10)\qquad f(t) = \frac{1}{3\sqrt{-3}}\,\frac{(t+1)(2t-1)(t-2)}{t(t-1)}
\]
et par suite
\[
\boxed{\ (g f)^*(t) = j = \frac{1}{-27}\Bigl(\frac{(t+1)(2t-1)(t-2)}{t(t-1)}\Bigr)^2\ }
\qquad (11)
\]
\note{devant $(gf)^*$, dans le cadre, un mot biffé illisible ; un trait
vertical sépare $(gf)^*(t)$ de $j$}
Cette formule, \ill{} \ill{} il \ill{} \uncertain{doit} \ill{}
\uncertain{définir} \ill{} \ill{} sur $J$ ! \ill{} $S$ \uncertain{lui-même},
et \ill{} \ill{}. Elle définit un revêtement \uncertain{galoisien} de groupe
$\mathfrak{S}_3$ de \struck{\ill{}} $(\mathbf{P}^1_S)$ \add{: fonction de
$(\mathbf{G}_m)_S$,} \ill{} \ill{} \ill{} non ramifié au-dessus de
\[
(\mathbf{G}_a)_S - (0) - (1) = (\mathbf{G}_m)_S - 1 .
\]
\note{l'indice du premier $\mathbf{G}$, lu $a$, est peu net ; la formule est
soulignée d'un double trait}

\uncertain{Comment} \ill{} \ill{} \ill{} les \uncertain{quotients} ?
\[
\begin{array}{ll}
\sigma & \mathbf{P}^1\\
& \big\downarrow f\\
\sigma' & \mathbf{P}^1\\
& \big\downarrow g\\
\sigma'' & \mathbf{P}^1
\end{array}
\]
\note{ce schéma est dans la marge gauche, en face du paragraphe qui suit}
On voit que $\sigma'(\infty)=0$, $\sigma'(0)=\infty$, $\sigma'(1)=1$,
\ill{} \ill{} \ill{} pour $\sigma''$, d'où \struck{$\sigma'(t)=1/t$}
\struck{et d'où \ill{} \ill{}}
\[
(12)\qquad \sigma'(t) = \sigma''(t) = 1/t
\]
et il serait facile de \ill{} \ill{} \uncertain{quotients} \ill{} \ill{}
\ill{} \ill{} \ill{} (p.\ ex.).
\note{la page s'arrête là, sur un double trait horizontal}

\end{document}
